\documentclass[11pt,a4paper]{article}

\usepackage[a4paper,margin=1.7cm]{geometry}
\usepackage[french]{babel}
\usepackage{enumitem}
\usepackage{xcolor}
\usepackage{array}

\definecolor{ink}{HTML}{181818}
\definecolor{muted}{HTML}{5C5C5C}
\definecolor{accent}{HTML}{C7922E}
\definecolor{panel}{HTML}{F5E8C7}
\definecolor{dark}{HTML}{141414}
\definecolor{soft}{HTML}{F7F3EA}

\setlength{\parindent}{0pt}
\setlength{\parskip}{0.65em}
\setlist[itemize]{leftmargin=1.1cm,itemsep=0.35em,topsep=0.4em}
\setlist[enumerate]{leftmargin=1.1cm,itemsep=0.45em,topsep=0.4em}

\newcommand{\guideTitle}{Apprendre les standards de jazz intelligemment}

\newcommand{\sectionbreakpage}{
  \newpage
}

\newcommand{\stepheader}[3]{
  {\fontsize{42}{46}\selectfont\bfseries\textcolor{accent}{#1}}
  \hspace{0.5em}
  \begin{minipage}[t]{0.8\linewidth}
    {\Huge\bfseries #2}

    {\Large\textcolor{muted}{#3}}
  \end{minipage}
  \vspace{0.8em}
}

\newcommand{\notebox}[2]{
  \vspace{0.6em}
  \noindent\fcolorbox{accent}{panel}{
    \parbox{\dimexpr\linewidth-2\fboxsep-2\fboxrule\relax}{
      {\large\bfseries\textcolor{accent}{#1}}

      #2
    }
  }
  \vspace{0.4em}
}

\begin{document}
\color{ink}

\begin{titlepage}
\pagecolor{dark}
\color{white}
\thispagestyle{empty}

\vspace*{1.5cm}
{\fontsize{21}{24}\selectfont\bfseries APPRENDRE}

\vspace{0.15cm}
{\fontsize{30}{33}\selectfont\bfseries LES STANDARDS}

\vspace{0.1cm}
{\fontsize{30}{33}\selectfont\bfseries DE JAZZ}

\vspace{0.1cm}
{\fontsize{21}{24}\selectfont\bfseries INTELLIGEMMENT}

\vfill

{\fontsize{34}{38}\selectfont\bfseries Le guide express}

\vspace{0.2cm}
{\fontsize{26}{30}\selectfont\bfseries en 5 etapes}

\vspace{0.5cm}
{\Large pour apprendre les morceaux}

\vspace{0.15cm}
{\Large de jazz comme un pro}

\vfill

{\Large Brent Vaartstra}

\vspace{0.8cm}
{\normalsize Version francaise adaptee a partir du guide fourni}
\end{titlepage}

\nopagecolor
\color{ink}
\setcounter{page}{1}

{\huge\bfseries Je m'appelle Brent Vaartstra}

\vspace{0.5em}

Je suis un musicien de jazz professionnel base a New York, auteur et podcasteur.
Mais on me connait surtout comme le musicien derriere le site internationalement reconnu
\textit{learnjazzstandards.com}. Au fil des annees, j'ai aide des milliers de musiciens qui
se heurtent exactement aux difficultes que nous allons aborder ici.

La bonne nouvelle, c'est que vous n'etes pas seul, et qu'il existe des solutions concretes
a tous ces problemes.

Je constate souvent que la racine du probleme se trouve dans le processus d'apprentissage.
Beaucoup de musiciens n'apprennent tout simplement pas les standards de jazz de la meilleure
maniere. C'est precisement pour cela que je vous propose une methode simple.

Dans ce guide, je vais vous montrer un processus en 5 etapes pour apprendre les standards
de jazz intelligemment. Si vous suivez ces etapes, vous connaitrez, jouerez et comprendrez
les standards bien mieux.

\notebox{Vous vous reconnaissez peut-etre ici}{
  \begin{itemize}
    \item Vous ne pouvez pas jouer un standard sans regarder la partition.
    \item Vous apprenez des standards, puis vous les oubliez rapidement.
    \item Vous vous perdez souvent dans la forme quand vous les jouez.
    \item Vous trouvez l'improvisation difficile parce que vous ne comprenez pas vraiment le morceau.
  \end{itemize}
}

Si vous avez deja ressenti l'une de ces frustrations, je comprends parfaitement. J'y suis
passe moi aussi, et cela peut etre tres decourageant.

\sectionbreakpage

{\huge\bfseries Apprendre a l'oreille vs. apprendre avec la partition}

\vspace{0.5em}

Quand il s'agit d'apprendre les standards de jazz correctement, la facon de s'y prendre
est cruciale. Le "comment" peut sembler secondaire, mais il ne l'est pas.

La solution la plus facile, c'est d'acheter une partition, de lire la melodie et les accords,
de memoriser tant bien que mal, puis de considerer que le travail est fait. C'est une voie
rapide et tres gratifiante sur le moment, mais ce n'est pas la meilleure maniere d'apprendre
le langage du jazz, ni la musique en general.

Souvenez-vous de ceci: le jazz est un langage. Et l'un des elements les plus importants dans
l'apprentissage d'un langage, c'est l'imitation. On entend quelque chose, puis on le reproduit.
N'importe quelle personne bilingue vous dira qu'elle est devenue fluide en s'entourant de
locuteurs natifs, en ecoutant, en dechiffrant, puis en repondant.

Le jazz n'a jamais ete une musique pensee d'abord pour etre apprise sur la page. Dans les annees
1940, a l'epoque du bebop, les musiciens se retrouvaient dans les clubs, s'ecoutaient les uns
les autres, et apprenaient sur le stand, en repetition, et en ecoutant les disques.

Cela ne signifie pas qu'ils ne savaient pas lire la musique. Mais apprendre a l'oreille restait
la methode principale. C'est un point essentiel: le jazz est avant tout une musique apprise
par l'ecoute. Si vous voulez devenir un excellent improvisateur, vous devez suivre cette trace.

Apprendre a l'oreille n'est pas toujours la solution la plus simple, surtout si vous n'y etes pas
habitue. Au debut, cela peut meme sembler difficile. Mais c'est de loin la meilleure approche,
et plus vous la pratiquez, plus elle devient naturelle.

La partition n'est pas l'ennemi pour autant. Elle peut jouer un role utile dans votre education
jazz. La lecture fait aussi partie de l'apprentissage d'un langage. Elle peut vous aider a
conceptualiser et a analyser ce que vous entendez.

Elle peut egalement servir a verifier le travail de votre oreille. Si vous avez appris un standard
a l'oreille, vous pouvez ensuite jeter un oeil a un accord incertain, ou a un fragment de melodie
qui vous semble ambigu.

Vous vous dites peut-etre: "D'accord, apprendre a l'oreille est plus efficace. Mais comment faire
concretement ?" C'est la qu'intervient mon processus LIST.

\sectionbreakpage

{\huge\bfseries Le processus LIST}

\vspace{0.5em}

Si vous etes comme moi, vous aimez les methodes claires, organisees, et etape par etape.
Une bonne checklist permet d'aller jusqu'au bout d'un apprentissage sans rien laisser de cote.

Voici donc mon petit acronyme pour apprendre les standards de jazz, et plus largement le langage
du jazz lui-meme: \textbf{LIST}.

\vspace{1.2em}

\begin{center}
\renewcommand{\arraystretch}{1.6}
\begin{tabular}{>{\centering\arraybackslash}m{0.2\linewidth} >{\centering\arraybackslash}m{0.2\linewidth} >{\centering\arraybackslash}m{0.2\linewidth} >{\centering\arraybackslash}m{0.2\linewidth}}
{\Huge\bfseries\textcolor{accent}{L}} & {\Huge\bfseries\textcolor{accent}{I}} & {\Huge\bfseries\textcolor{accent}{S}} & {\Huge\bfseries\textcolor{accent}{T}} \\
{\large\bfseries Listen} & {\large\bfseries Internalize} & {\large\bfseries Sing} & {\large\bfseries Transfer} \\
{\normalsize Ecouter} & {\normalsize Interioriser} & {\normalsize Chanter} & {\normalsize Transferer}
\end{tabular}
\end{center}

\vspace{1em}

Cette sequence vous oblige a commencer par l'ecoute, a faire entrer la musique dans votre oreille,
puis dans votre voix, avant de la faire passer sur votre instrument. C'est ce qui rend
l'apprentissage beaucoup plus solide.

\sectionbreakpage

\stepheader{L}{Listen}{Ecouter}

La premiere etape est tres simple: ecoutez le standard. Cela peut sembler evident, mais beaucoup
d'etudiants se precipitent sur leur instrument alors qu'ils ont a peine ecoute le morceau.

La chose la plus importante a faire lorsque vous apprenez un standard, c'est de poser votre
instrument et d'ecouter. Trouvez autant d'enregistrements que possible du morceau que vous
voulez apprendre et parcourez-les. Vous devez d'abord faire connaissance avec la chanson.
Si vous sautez cette phase, vous partez deja du mauvais pied.

Ce n'est qu'apres avoir repere plusieurs versions et commence a vraiment les ecouter que vous
devez passer a l'etape suivante.

\notebox{ASTUCE PRO \#1 - Apprenez les paroles}{
  Apprendre les paroles aide enormement a interioriser une melodie. Dans le cas de nombreux
  standards de jazz, les textes sont poetiques, memorables, et racontent l'essence meme du morceau.
}

\sectionbreakpage

\stepheader{I}{Internalize}{Interioriser}

Cette etape implique encore plus d'ecoute, mais une ecoute differente: une ecoute intentionnelle.

Imaginez que vous soyez assis dans votre salon devant un film. Si le film vous captive, vous pouvez
rester concentre pendant deux heures ou plus. C'est impressionnant quand on y pense.

Et si vous traitiez la musique de la meme maniere ? En general, il est deja excellent d'ecouter
la musique en lui accordant toute votre attention, sans distraction. Mais c'est particulierement
important lorsque vous essayez d'apprendre un nouveau standard.

Si vous lui consacrez une attention reelle, le morceau commencera a s'interioriser. Il s'ancrera
petit a petit dans votre subconscient, et il vous deviendra familier bien avant que vous le jouiez.

\sectionbreakpage

\stepheader{S}{Sing}{Chanter}

Cette etape est vraiment importante. Chanter est un moyen tres puissant de verifier que vous avez
reellement interiorise ce que vous entendez. Non, vous n'avez pas besoin d'etre un grand chanteur,
et oui, vous pouvez aussi fredonner ou siffler si vous preferez.

L'application principale du chant concerne la melodie. Assurez-vous de pouvoir chanter la melodie
du standard, a la fois avec l'enregistrement et seul, avant de la travailler sur votre instrument.

Le chant retire en quelque sorte 50\% du travail. Il prouve que vous avez deja integre l'information.
Il ne reste plus qu'a la transferer vers votre instrument.

Si vous voulez aller plus loin, essayez de chanter aussi les notes de basse. Cela vous aidera a
commencer a entendre l'harmonie et la progression d'accords du standard. Ecoutez la basse sur vos
enregistrements et voyez si vous pouvez chanter les fondamentales. C'est une astuce plus avancee,
mais elle vous aidera enormement pour l'etape suivante.

\notebox{ASTUCE PRO \#2 - Trouver la version de base de la melodie}{
  Une remarque qui revient souvent est la suivante: "Quelle est la bonne melodie, si les jazzmen
  prennent tant de libertes ?" Souvenez-vous de ceci: Frank Sinatra est votre ami. Il chante tres
  souvent les melodies de facon droite. Si le morceau est moins classique, cherchez l'artiste ou
  le compositeur qui l'a interprete en premier. Il y a de fortes chances qu'il joue la melodie
  au plus proche de son intention d'origine.
}

\sectionbreakpage

\stepheader{T}{Transfer}{Transferer}

Il est maintenant temps de reprendre votre instrument. Jusqu'ici, vous n'auriez pas du vraiment
y toucher. A ce stade, vous devriez deja bien connaitre le morceau, connaitre la melodie, et
avoir au moins une premiere idee des accords selon votre niveau actuel.

\begin{enumerate}
  \item \textbf{Transferer la melodie.} Commencez a jouer la melodie sur votre instrument. Cela
  devrait surtout consister a retrouver sous vos doigts ce que vous savez deja chanter.
  \item \textbf{Apprendre les accords.} C'est la partie la plus exigeante lorsqu'on apprend a
  l'oreille, mais faites de votre mieux. Il est extremement utile d'avoir une bonne comprehension
  du fonctionnement de l'harmonie jazz. L'idee est de combiner votre connaissance des progressions
  d'accords avec votre oreille afin d'identifier la qualite des accords - majeur, mineur,
  dominant, etc. - et leur fonction dans la progression.
\end{enumerate}

\notebox{ASTUCE PRO \#3 - Ecoutez les notes de basse}{
  Comme dans l'etape "Chanter", le mouvement de basse est capital. Si vous identifiez par exemple
  que le bassiste joue un Mi bemol de concert sur le premier accord, il ne vous reste qu'a trouver
  la qualite de cet accord. Puis vous ecoutez la basse de l'accord suivant. Si la note devient
  un Do de concert, vous pouvez deja commencer a reconnaitre le mouvement harmonique.

  Si l'apprentissage des progressions a l'oreille est tout nouveau pour vous et que vous bloquez,
  ce n'est pas grave. C'est alors un bon moment pour regarder la partition. Mais faites-le en
  l'etudiant avec l'enregistrement, en arretant et relancant si necessaire, afin de continuer a
  relier ce que vous voyez a ce que vous entendez.
}

\sectionbreakpage

\stepheader{S}{Study}{Etudier}

J'aime souvent ajouter une lettre supplementaire au processus LIST et le transformer en \textbf{LISTS}.
Si vous voulez vraiment entrer dans le morceau et le comprendre, vous devez l'etudier.

Analysez les progressions d'accords pour voir ce qui s'y passe reellement, puis utilisez des outils
concrets du jazz pour commencer a improviser dessus.

L'analyse des standards avec les degres en chiffres romains peut vous aider a comprendre la fonction
des accords dans une progression. Lorsque vous commencez a decomposer les choses et a les
conceptualiser, tout devient plus clair. Si vous comprenez la theorie, improviser devient beaucoup
plus naturel.

Ensuite, cartographiez les notes d'accord, les guide tones, et les gammes utiles. Cela vous aidera
a identifier des choix de notes pertinents pour votre improvisation, au lieu d'avoir l'impression
d'avancer au hasard.

Enfin, apprenez un etude ou un solo de jazz pour assimiler du "langage jazz", autrement dit pour
imiter la maniere dont les grands musiciens "parlent". Puis allez jouer, improvisez, et mettez
tout cela en pratique.

\vspace{0.8em}
\rule{\linewidth}{0.6pt}
\vspace{0.8em}

{\Large\bfseries A propos de l'auteur}

\vspace{0.35em}

Brent Vaartstra est un guitariste de jazz professionnel et un pedagogue base a New York.
Il est le createur du blog et du podcast \textit{learnjazzstandards.com}, dont la mission est de
rendre l'apprentissage du jazz plus simple et plus plaisant pour tous les instrumentistes, quel
que soit leur niveau.

\vspace{1em}

{\Large\bfseries Avant de partir...}

\vspace{0.35em}

Pensez a jeter un oeil au programme \textit{LJS Inner Circle}. Plus de 1\,000 musiciens, sur des
instruments tres varies, s'y retrouvent pour faire progresser leur jeu jazz de maniere concrete.

Vous pouvez aussi vous abonner au \textit{Learn Jazz Standards Podcast} sur votre application
preferee. Un nouvel episode sort chaque semaine avec des lecons, conseils et strategies pour
vous aider a devenir un meilleur musicien de jazz, quel que soit votre instrument.

\end{document}
